\begin{frame}{자주 묻는 질문 (2/2)}
  \begin{itemize}
    \item \kq{Q. ``Temporal difference''가 실제로 차이를 이루는 양은?}
      $\delta = r + \gamma V(s') - V(s)$ 안에서 $V(s') - V(s)$가
      ``$t{+}1$ 추정치와 $t$ 추정치의 차''. 관찰된 $r$에 더해
      \textbf{인접한 두 시점의 가치 추정치의 차}가 학습 신호를 만든다.
      MC(리턴과 추정치의 차)와 DP(정확한 기댓값과 추정치의 차)는 이
      ``차''의 한쪽을 추정치가 아닌 것으로 채우므로 ``temporal
      difference''라 부르지 않는다.
    \item \kq{Q. 정책이 하나뿐이라 ``예측''만 하고 ``개선''은 안 하나?}
      그렇다 --- 이번 절(5.1절 MC 예측과 함께)은 \textbf{정책 예측} 문제.
      ``더 나은 정책으로 갈아타기''는 6.2절에서 $V(s) \to Q(s,a)$로 바꾸고,
      행동 선택이 $\arg\max_a Q(s,a)$로 읽히는 순간부터.
      TD 갱신식 뼈대는 \textbf{전혀} 바뀌지 않는다 --- 바뀌는 것은
      ``어떤 가치함수를 갱신하느냐''뿐.
  \end{itemize}
\end{frame}
